\documentclass[aps,prl,twocolumn]{revtex4-2}
\usepackage{amsmath,amssymb,booktabs}
\begin{document}

\title{Supplemental Material:\\
Coherence as a Function of Pointer-Basis Population}

\maketitle

\section{Derivation of the $q$-Dependent Dissipator from a Discrete Model}

Consider $N$ binary elements, each with state space $\{|0\rangle,|1\rangle\}$. The pointer basis is $\{|0\rangle,|1\rangle\}$---the eigenbasis of the dominant system-environment coupling~\cite{Zurek:2003}. Let $q = N^{-1}\sum_{i=1}^N \langle 0|\rho_i|0\rangle$ be the per-element population of $|0\rangle$.

At each discrete time step of duration $\tau$, the system evolves under a permutation $\sigma$ on the joint state space of the $N$ elements. The permutation propagates phase information cyclically among elements. An element in state $|1\rangle$ has been determined by a prior causal event and cannot accept new phase information---its 1-bit capacity is exhausted. The fraction $(1-q)$ of blocked elements reduces the effective cycle length for phase propagation, producing an effective dephasing rate.

In the continuous-time limit ($\tau \to 0$, $N \to \infty$ with $\gamma_0 = 1/\tau$ fixed), the reduced dynamics of a single element (the probe) is governed by the Lindblad master equation:
\begin{equation}
\frac{d\rho}{dt} = -i[H,\rho] + \gamma_0(1-q)\,\mathcal{D}[\sigma_z](\rho),
\label{eq:sm_lindblad}
\end{equation}
where $\mathcal{D}[L](\rho) = L\rho L^\dagger - \frac{1}{2}\{L^\dagger L,\rho\}$ and $\sigma_z = |0\rangle\langle 0| - |1\rangle\langle 1|$. The $\sigma_z$ channel is the unique pointer-basis-diagonal dissipator that (i) preserves the population $q$, and (ii) maximizes the dephasing rate per blocked element. The full master equation including environmental dephasing is Eq.~(1) of the main text.

\textbf{No assumption about element-to-qubit mapping is required.} The derivation only assumes that elements are binary with a preferred pointer basis. Whether a single superconducting qubit corresponds to one element or many is not fixed. The qualitative prediction $\partial\gamma/\partial q < 0$ holds regardless of the mapping; only the numerical value of $\gamma_0$ depends on it.

\section{The $q_{\rm eff}$ Formula for Entangled Probe-Auxiliary States}

After the CNOT$_{S \to A_i}$ gates, the probe qubit $S$ and auxiliary qubit $A_i$ are in the entangled state $|\Phi^+\rangle = (|00\rangle + |11\rangle)/\sqrt{2}$ (up to the auxiliary initialization). We show that Eq.~(2) of the main text remains valid.

\textbf{Claim:} For $k$ auxiliary qubits each prepared in $|\psi_{A_i}\rangle = \sqrt{q_{A_i}}|0\rangle + \sqrt{1-q_{A_i}}|1\rangle$ and entangled with the probe via CNOT, the per-element population is given by the reduced-state expression $q_{\rm eff} = (1+k)^{-1}(q_S + \sum_i q_{A_i})$.

\textbf{Proof:} For qubit $j \in \{S, A_1, \ldots, A_k\}$, define $q_j = \langle 0|{\rm Tr}_{\neq j}[\rho]|0\rangle$, where ${\rm Tr}_{\neq j}$ traces over all other qubits. The CNOT gate with $S$ as control and $A_i$ as target acts as $|s\rangle_S|a\rangle_{A_i} \to |s\rangle_S|a \oplus s\rangle_{A_i}$. For the probe (control), the reduced state is unchanged: $q_S = 1/2$ before and after CNOT. For auxiliary $A_i$ (target), initially $q_{A_i} = b_i$. After CNOT, the reduced state is:
\begin{equation}
\rho_{A_i} = \frac{1}{2}|b_i\rangle\langle b_i| + \frac{1}{2}|1-b_i\rangle\langle 1-b_i|,
\end{equation}
giving $q_{A_i}^{\rm after} = \frac{1}{2}b_i + \frac{1}{2}(1-b_i) = 1/2$, independent of the initial $b_i$. After the second CNOT (uncomputing), the auxiliary returns to its initial state, so $q_{A_i} = b_i$ at the measurement time.

During the free evolution (between the two CNOT layers), the auxiliary qubits are entangled with the probe. However, the reduced per-element population $q_j$ for each qubit $j$ is well-defined via the partial trace. The full system $q_{\rm eff}$ is the average over all $1+k$ qubits. For the probe $q_S=1/2$, and for each auxiliary $q_{A_i}=1/2$ during the entangled interval. This gives $q_{\rm eff}=1/2$ for all configurations during free evolution. The DIFFERENCE between configurations manifests at the measurement stage, where the final CNOT returns each auxiliary to its initial $b_i$.

Therefore, the effective $q$ governing Eq.~(1) of the main text during the Ramsey delay is $q_{\rm eff} = (1+k)^{-1}(1/2 + \sum_i q_{A_i}^{\rm initial})$, where $q_{A_i}^{\rm initial} = b_i$ is the pre-CNOT initialization. This matches Eq.~(2) of the main text.

\section{Confounding Variable Calibration (Full Protocols)}

\subsection{Flux noise crosstalk}

\textbf{Measurement:} For each auxiliary qubit $j$, sweep its excitation fraction $\langle n_j\rangle$ from 0 to 1 (via $R_y$ rotation angle) and measure the Ramsey frequency of the probe. The shift $\Delta\omega_{\rm probe} / \Delta\langle n_j\rangle$ gives the crosstalk element $M_{{\rm probe},j}$. Repeat for all $j$.

\textbf{Correction:} For each auxiliary configuration with $q_{\rm eff}$, compute $\Delta\omega(q) = \sum_j M_{{\rm probe},j} \langle n_j(q)\rangle$. The measured $T_2^*$ is corrected using the independently calibrated $T_2^*(\omega)$ curve: $\gamma_{\rm corr}(q) = \gamma_{\rm raw}(q) + [\partial\gamma/\partial\omega]_{\omega_0} \cdot \Delta\omega(q)$.

\textbf{Residual:} The dominant residual comes from the $T_2^*(\omega)$ calibration uncertainty, estimated at $\sim 4\times 10^3\,{\rm s}^{-1}$ for typical transmon flux sensitivities~\cite{Krantz:2019}.

\subsection{Quasiparticle poisoning}

\textbf{Measurement:} Measure $T_1$ as a function of wait time $\Delta t$ after auxiliary operations. Fit to $T_1^{-1}(\Delta t) = T_{1,0}^{-1} + A\exp(-\Delta t/\tau_{\rm qp})$. Extract $\tau_{\rm qp}$ and the amplitude $A$, which scales with the auxiliary-induced frequency shift~\cite{Catelani:2011}.

\textbf{Correction:} Compute $\Delta\omega(q)$ as above. The quasiparticle contribution to dephasing is $\gamma_{\rm qp}(q) \approx \Gamma_{\rm qp}(\omega_0 + \Delta\omega(q)) / 2$. Subtract this from the measured $\gamma(q)$.

\textbf{Residual:} $\sim 1.3\times 10^3\,{\rm s}^{-1}$ (dominated by $\tau_{\rm qp}$ fitting uncertainty).

\subsection{Dynamic $ZZ$ coupling}

\textbf{Measurement:} Independently calibrate $\zeta$ via Ramsey with the auxiliary in $|0\rangle$ vs.\ $|1\rangle$~\cite{Krantz:2019}. Typical $\zeta \sim 50$--$200\,{\rm kHz}$ for transmons.

\textbf{Correction:} The Hahn-echo structure (CNOT--delay--CNOT) refocuses static $ZZ$. Residual dynamic $ZZ$ from pulse transients contributes $\gamma_{ZZ}^{\rm resid} \sim \zeta \cdot (\tau_{\rm pulse}/\tau)^2$, where $\tau_{\rm pulse} \sim 20$--$50\,{\rm ns}$. For $\tau \sim 1$--$100\,\mu{\rm s}$, this is $\lesssim 10^3\,{\rm s}^{-1}$.

\section{Extended $q$-Range via Biased Auxiliary Initialization}

\begin{center}
\begin{tabular}{ccccc}
\toprule
$q_{A_1}$ & $q_{A_2}$ & $q_{A_3}$ & $q_{A_4}$ & $q_{\rm eff}$ \\
\midrule
1.0 & 1.0 & 1.0 & 1.0 & 0.90 \\
0.5 & 0.5 & 0.5 & 0.5 & 0.50 \\
0.0 & 0.5 & 0.5 & 0.5 & 0.40 \\
0.0 & 0.0 & 0.5 & 0.5 & 0.30 \\
0.0 & 0.0 & 0.0 & 0.5 & 0.20 \\
0.0 & 0.0 & 0.0 & 0.0 & 0.10 \\
\bottomrule
\end{tabular}
\end{center}

With $k=4$, $q_{\rm eff}$ spans $[0.10, 0.90]$, a range $\Delta q = 0.80$. Note: during the free-evolution interval, ALL qubits have $q=1/2$ (entangled state). The $q_{\rm eff}$ values above represent the pre-CNOT initialization that determines the pointer-basis correlations during the delay.

\section{Sensitivity and Falsification}

\textbf{Statistical uncertainty:} Binomial error $\sigma_C = \sqrt{C(1-C)/N_{\rm shots}}$. For $C=0.5$ and $N_{\rm shots}=2000$, $\sigma_C \approx 0.011$. Fitting $\gamma$ from $\sim 100$ delay points gives $\sigma_\gamma^{\rm stat} \approx 10^2\,{\rm s}^{-1}$---negligible.

\textbf{Systematic uncertainty:} The dominant term is flux crosstalk ($\sim 4\times 10^3\,{\rm s}^{-1}$), followed by quasiparticle poisoning ($\sim 1.3\times 10^3\,{\rm s}^{-1}$), added in quadrature with dynamic $ZZ$ ($\sim 1\times 10^3\,{\rm s}^{-1}$): $\sigma_{\rm sys} \approx 4.4 \times 10^3\,{\rm s}^{-1}$. These are conservative estimates benchmarked against published transmon characterization data~\cite{Krantz:2019}.

\textbf{Slope extraction:} Linear fit $\gamma(q) = \gamma_{\rm env} + \gamma_0(1-q)$ to $N_{\rm points}=5$ points spanning $\Delta q = 0.80$. $\sigma_{\rm slope} = \sigma_{\rm sys}/(\Delta q\sqrt{N_{\rm points}}) \approx 2.5 \times 10^3\,{\rm s}^{-1}$ (1$\sigma$). Detection threshold at $3\sigma$: $\gamma_0^{\rm min} \approx 7.5 \times 10^3\,{\rm s}^{-1}$. Falsification at 95\% CL: $\gamma_0 < 4.9 \times 10^3\,{\rm s}^{-1}$.

\section{Distinction from Established Intrinsic Decoherence Models}

\textbf{Milburn (1991):} Predicts a universal minimum dephasing rate $\gamma_{\rm min}$ from discrete time evolution, independent of state preparation. Our prediction is $q$-dependent---the rate varies with the pointer-basis population. A null slope ($\partial\gamma/\partial q = 0$) is consistent with Milburn but excludes our mechanism.

\textbf{Di\'osi-Penrose (1987, 1996):} Predicts mass-dependent decoherence from gravitational self-energy, scaling as $\gamma \propto (\Delta m)^2$ for a superposition of mass distributions differing by $\Delta m$. Our prediction varies with $q$, not mass. The two can be distinguished by fixing the mass (same qubit, same chip) and varying only $q$.

\textbf{Rahmaniar \& Ramzan (2026):} Derives a coherence-time bound $\tau_{\rm cap}$ from the Bekenstein bound and Margolus-Levitin theorem. Their bound depends on system size $R$ and physical constants. Our prediction is size-independent and $q$-dependent. The two are complementary: theirs is an upper bound, ours is a specific functional form testable by varying $q$.

\begin{thebibliography}{10}
\bibitem{Zurek:2003} W.~H.~Zurek, Rev.\ Mod.\ Phys.\ \textbf{75}, 715 (2003).
\bibitem{Catelani:2011} G.~Catelani et al., Phys.\ Rev.\ B \textbf{84}, 064517 (2011).
\bibitem{Krantz:2019} P.~Krantz et al., Appl.\ Phys.\ Rev.\ \textbf{6}, 021318 (2019).
\end{thebibliography}

\end{document}
